%% Dans ce fichier, les �ventuels accents sont cod�s
%% en ISO-8859-15 (aka latin9)
%% La ligne suivante donne trois exemples 
%% e accent aigu : �, e accent grave : �, e-dans-l'o : �
%%%%%ENTREES%%%%
\begin{tabarg}
\multicolumn{3}{l}{\titretab {Entr\'ees}}\\
\hline
$X$ & Vecteur, \texttt{float[$N_{level}$]} & Signal lidar \\
$Z$ & Vecteur, \texttt{float[$N_{level}$]} & Altitudes r�elles des portes du lidar (m)\\
$N_{level}$ & Coeff., \texttt{float} & Nombre de portes (\$Arg1.array\_dim)\\
$Z_1$ & Coeff., \texttt{float} & Altitude minimale par rapport au lidar de la zone de recherche (m)\\
$Z_2$ & Coeff., \texttt{float} & Altitude maximale par rapport au lidar de la zone de recherche (m)\\
$N_{step}$ & Coeff., \texttt{int} & Nombre de points de la moyenne glissante de $X$ \\
\end{tabarg}



\medskip
%%%%%SORTIES%%%%
\begin{tabarg}
\multicolumn{3}{l}{\titretab {Sorties}}\\
\hline
$Z_{ref}$ & Vecteur, \texttt{float} & Altitude de r�f�rence du signal mol�culaire  (m)\\
\end{tabarg}
